\subsection{The classical scalar comparison}
\label{sec:v15-scalar}
Smooth linearization on a one-dimensional contracting branch is
classical. It requires neither a symplectic chart on a neighborhood
nor an analytic hypothesis. This is the one-dimensional case of
Sternberg's theorem; see~\cite{Sternberg} and the precise regularity
statement in the introduction of Eynard-Bontemps--Navas~\cite{EN}.
We record the elementary product construction in the normalization
needed here. The scalar comparison and its application to the width
were pointed out in the memorandum~\cite{A2ReviewV14}. Its role is to
separate a classical coordinate from the identification of the physical
determinant amplitude.

\begin{lemma}[Normalized scalar linearization]
\label{lem:v15-scalar-product}
Let $R_\xi$ be a jointly smooth family on a fixed interval about zero,
with parameter $\xi$ in a compact set, such that
\[
 R_\xi(0)=0,\qquad R_\xi'(0)=\lambda_\xi,\qquad
 0<\ell\le R_\xi'(u)\le q<1.
\]
The family is understood to be defined on an open neighborhood of the
compact parameter set. On a common smaller interval the series
\begin{equation}\label{eq:v15-scalar-product}
 \begin{split}
 h_\xi(u)&=\log\frac{R_\xi'(u)}{\lambda_\xi},\\
 L_\xi(u)&=\sum_{i=0}^{\infty}h_\xi(R_\xi^i(u)),\qquad
 \mathfrak B_\xi(u)=\exp L_\xi(u),\qquad
 \mathfrak z_\xi(u)=\int_0^u\mathfrak B_\xi(x)\,\dd x
 \end{split}
\end{equation}
converges with every fixed mixed parameter and endpoint derivative.
The coordinate $\mathfrak z_\xi$ is increasing, normalized by
$\mathfrak z_\xi(0)=0$ and $\mathfrak z_\xi'(0)=1$, and satisfies
\begin{equation}\label{eq:v15-scalar-cocycle}
 \mathfrak B_\xi(R_\xi(u))R_\xi'(u)
       =\lambda_\xi\mathfrak B_\xi(u),\qquad
 \mathfrak z_\xi\circ R_\xi=\lambda_\xi\mathfrak z_\xi.
\end{equation}
It is the unique normalized differentiable local linearizer. For every
fixed $k$ and $\sigma\in(\lambda_+,1)$, where
$\lambda_+=\max_\xi\lambda_\xi$, one has
\begin{equation}\label{eq:v15-scalar-iterates}
 \|\lambda_\xi^{-n}R_\xi^n-\mathfrak z_\xi\|_{C^k}
 +\|\lambda_\xi^{-n}(R_\xi^n)'-\mathfrak B_\xi\|_{C^k}
                                      \le C_{k,\sigma}\sigma^n.
\end{equation}
Here and below such norms include any fixed collection of mixed
parameter and endpoint derivatives on the common domain.
\end{lemma}
\begin{proof}
Suppress $\xi$ in compositions, but not in differentiation. Put
$x_n=R_\xi^n(u)$. Since $R_\xi(0)=0$ and $|R_\xi'|\le q$,
$|x_n|\le q^n|u|$. At every fixed positive mixed order $m$,
differentiation of $x_{n+1}=R_\xi(x_n)$ leaves
$R_\xi'(x_n)$ multiplying the highest derivative of $x_n$.
Every other term either contains a lower derivative of $x_n$, or
is a pure parameter derivative of $R_\xi$ evaluated at $x_n$.
The latter vanishes at $x_n=0$, because the fixed point is common to
the family. Induction and the mean value theorem therefore give
\[
 \sup_{|\alpha|+r=m}|\partial_\xi^\alpha\partial_u^r x_n|
                  \le C_m(1+n)^{M_m}q^n.
\]
Indeed, at the inductive step the right-hand forcing in the highest
order recurrence is bounded by a polynomial in $n$ times $q^n$;
iteration of that recurrence only increases the polynomial degree.
Uniform bounds on the finitely many derivatives of $R_\xi$ involved
follow from compactness. This argument covers endpoint derivatives,
parameter derivatives and their mixtures.

The function $h_\xi$ is smooth, $h_\xi(0)=0$, and all its pure
parameter derivatives vanish at zero. The same chain-rule argument
bounds each fixed mixed derivative of $h_\xi(x_n)$ by
$C_m(1+n)^{N_m}q^n$. Summability proves convergence of
\eqref{eq:v15-scalar-product} with all the asserted derivatives.
In particular $L_\xi(0)=0$ and $\mathfrak B_\xi(0)=1$.
Shifting the summation index gives
$L_\xi(u)-L_\xi(R_\xi(u))=h_\xi(u)$. Exponentiation proves the
first identity in \eqref{eq:v15-scalar-cocycle}. Integration from
zero proves the second, since both sides vanish there.

If $\widetilde{\mathfrak z}$ is another normalized differentiable
linearizer, set $H=\widetilde{\mathfrak z}\circ\mathfrak z^{-1}$.
Then $H(0)=0$, $H'(0)=1$, and $H(\lambda z)=\lambda H(z)$.
For every $z$ in a sufficiently small interval,
$H(z)=\lambda^{-n}H(\lambda^nz)\to z$. This proves uniqueness.

The exact relation
$R_\xi^n=\mathfrak z_\xi^{-1}\circ(\lambda_\xi^n\mathfrak z_\xi)$
gives the sharper rate based on $\lambda_+$ rather than the
preliminary contraction bound $q$. The inverse function theorem on
a common smaller interval makes $\mathfrak z_\xi^{-1}$ jointly
smooth. Its value and first derivative at zero are $0$ and $1$;
write it as $x+x^2A_\xi(x)$. Thus
\[
 \lambda_\xi^{-n}R_\xi^n(u)-\mathfrak z_\xi(u)
 =\lambda_\xi^n\mathfrak z_\xi(u)^2
      A_\xi(\lambda_\xi^n\mathfrak z_\xi(u)).
\]
Every fixed mixed derivative is bounded by
$C_k(1+n)^{N_k}\lambda_+^n$; the positive lower bound on
$\lambda_\xi$ controls derivatives of its powers. Apply the same
argument with one additional endpoint derivative for the density
estimate. A fixed strict margin $\lambda_+<\sigma$ absorbs the
polynomial factors. Finally, the finite product identity
\begin{equation}\label{eq:v15-finite-product}
 \exp\!\left\{\sum_{i=0}^{n-1}h_\xi(R_\xi^i(u))\right\}
                   =\lambda_\xi^{-n}(R_\xi^n)'(u)
\end{equation}
follows directly from the chain rule and agrees with the density limit.
\end{proof}

The alternating scalar identity is equally classical. Suppose that
$\varphi_b(0)=0$, $\varphi_b'(0)=r_b>0$ and
$R_b=\varphi_{1-b}\circ\varphi_b$, with $r_0r_1=\lambda<1$.
Let $\mathfrak z_b$ be the normalized scalar linearizer of $R_b$.
The identity $\varphi_b\circ R_b=R_{1-b}\circ\varphi_b$ shows that
$r_b^{-1}\mathfrak z_{1-b}\circ\varphi_b$ is another normalized
linearizer of $R_b$. Uniqueness gives
\[
 \mathfrak z_{1-b}\circ\varphi_b=r_b\mathfrak z_b,\qquad
 \mathfrak B_{1-b}(\varphi_b(u))\varphi_b'(u)
                                  =r_b\mathfrak B_b(u).
\]
These scalar facts by themselves do not identify the amplitude in a
normalized finite physical bridge. That identification is the content
of the following theorem and its exact Schur-concatenation proof.
